\documentclass[aspectratio=169]{beamer}

\usepackage[spanish]{babel}
\usepackage[utf8]{inputenc}
\usepackage[T1]{fontenc}
\usepackage{lmodern}
\usepackage{amsmath, amsfonts}
\usepackage{hyperref}

\usetheme{Madrid}

\title{Complejidad computacional}
\author{Sebastian Mestre}
\date{}

\begin{document}

\begin{frame}
  \titlepage
\end{frame}


\begin{frame}
  \centering
  \Huge
  \textbf{complejidad $\neq$ costo}
\end{frame}

\begin{frame}{Costo vs complejidad}
  \begin{block}{\textbf{Definición: complejidad}}
      La complejidad computacional \textbf{de un problema} es el mínimo costo asintótico posible que tiene \textbf{un programa} que resuelve ese problema.

  \end{block}
    
  \begin{itemize}
  
    \item Ejemplo: para ordenar un arreglo de elementos que solo podemos comparar de a pares, es sencillo diseñar un algoritmo de costo $O(N^2)$.
    \item Existen algoritmos más sofisticados (por ejemplo, \emph{merge sort}, \emph{quick sort}) con costo $O(N \log N)$.
    \item Est\'a demostrado que \textbf{no existen} algoritmos asintóticamente más rápidos.
    \item Entonces, el problema de ordenar un arreglo tiene complejidad $O(N \log N)$.
  \end{itemize}
\end{frame}

\section{P y NP}

\begin{frame}{La clase P}
  \begin{block}{\textbf{Definición: P}}
    $P$ es el conjunto de todos los problemas cuya complejidad computacional es \textbf{polinomial}.

    Es decir, problemas que se pueden resolver con costo:
    \[
      O(N),\; O(N^2),\; O(N^3),\; O(1),\; O(N \log N), \dots
    \]
    
    Decimos que estos problemas tienen \textbf{complejidad polinomial}.
  \end{block}
    
\end{frame}

\begin{frame}{Problemas fuera de P}
  \begin{itemize}
    \item Si un problema no es $P$, tiene complejidad \textbf{mayor que cualquier polinomio}.
    \item Ejemplos típicos de costos:
      \[
        O(2^N),\quad O(N!),\quad \text{o peores.}
      \]
    \item Una forma de pensarlos: \textbf{“problemas donde no te queda más opción que probar todas las posibilidades”}.
  \end{itemize}
\end{frame}

\begin{frame}{La clase NP}
  \begin{block}{\textbf{Definición: NP}}
    $NP$ es el conjunto de todos los problemas para los cuales verificar si una solución dada es correcta tiene \textbf{complejidad polinomial}.
  \end{block}
  \begin{itemize}
    \item En la práctica, verificar una solución suele ser más fácil que \textbf{encontrarla}.
    \item La teoría respalda parcialmente esta idea: $P \subseteq NP$.
    \item En otras palabras: si un problema se puede \textbf{resolver} en tiempo polinomial,
      entonces también se puede \textbf{verificar} una solución en tiempo polinomial.
  \end{itemize}
\end{frame}

\section{Reducciones}


\begin{frame}{Reducciones}
    \begin{itemize}
        \item En programación competitiva usamos todo el tiempo algoritmos conocidos para resolver problemas nuevos.
        \item En realidad, esos algoritmos son soluciones a problemas conocidos, y nosotros \textbf{reducimos} nuestro problema a uno de esos problemas conocidos.
        \begin{block}{\textbf{Definición: reducción}}
          Una reducción es una funcion que convierte un caso de prueba de un problema a un caso de prueba de otro problema que tiene la misma respuesta.
        \end{block}
    \end{itemize}
\end{frame}

\section{NP-completos}

\begin{frame}{Problemas NP-completos}
    \begin{block}{\textbf{Definición: NP-completo}}
      Un problema es NP-completo si es NP y cualquier otro problema de NP se puede reducir a él con costo de reducción polinomial.
      
      Es decir, podemos tomar cualquier caso de prueba de cualquier problema de $NP$ y, en tiempo polinomial, convertirlo en un caso de prueba de uno de estos problemas especiales.
    \end{block}

    \begin{itemize}
        \item \href{https://en.wikipedia.org/wiki/Boolean_satisfiability_problem}{\textbf{SAT} (satisfacibilidad booleana).}
        \item \href{https://en.wikipedia.org/wiki/Travelling_salesman_problem}{\textbf{TSP} (problema del viajante).}
        \item \textbf{Max-clique}.
        \item \textbf{Subset-sum}.
        \item \textbf{Max-independent set}.
        \item \textbf{Set-cover}.
  \end{itemize}
\end{frame}

\section{P vs NP}

\begin{frame}{La gran pregunta: P vs NP}
  \begin{itemize}
    \item Sabemos que $P \subseteq NP$.
    \item Pero nadie sabe si existe algún problema en $NP$ que \textbf{no} esté en $P$.
    \item En particular, se cree que los problemas \textbf{NP-completos} no pertenecen a $P$.
  \end{itemize}
\end{frame}

\begin{frame}{Si un NP-completo estuviera en P...}
  \begin{itemize}
    \item Si encontráramos una \textbf{solución polinomial} para algún problema NP-completo:
      \begin{itemize}
        \item Podríamos reducir \textbf{cualquier} problema de $NP$ a ese problema NP-completo.
        \item Resolveríamos también esos problemas en tiempo polinomial.
      \end{itemize}
    \item Por lo tanto, se tendría que:
      \[
        P = NP.
      \]
  \end{itemize}
\end{frame}

\section{NP-completo en programación competitiva}

\begin{frame}{NP-completo en programación competitiva}
  \begin{itemize}
    \item Un NP-completo \textbf{no se puede resolver} rapido para $N$ grande. Sin embargo, muchos \textbf{casos especiales} salen con técnicas estándar.
    \item Clave: si el problema se parece a uno NP-completo famoso pero tiene una \textbf{pequeña diferencia o restricción}, normalmente:
      \begin{itemize}
        \item Esa diferencia es lo que hay que explotar para resolverlo eficientemente.
      \end{itemize}
    \item A veces se toma un problema NP-completo con $N$ chico. No olvidemos practicar \textbf{búsquedas exhaustivas}:
      \begin{itemize}
        \item DP con bitmasks
        \item \texttt{next\_permutation}
        \item backtracking optimizado
      \end{itemize}
  \end{itemize}
\end{frame}

\section*{Para leer más}

\begin{frame}{Para leer más}
  \begin{itemize}
    \item Material:
      \begin{itemize}
        \item \href{https://oia-politecnico.github.io/wiki/complejidad}{Wiki de complejidad (OIA Politécnico)}
        \item \href{https://en.wikipedia.org/wiki/NP-complete}{NP-complete (wikipedia)}
        \item \href{https://en.wikipedia.org/wiki/Boolean_satisfiability_problem}{SAT (wikipedia)}
        \item \href{https://en.wikipedia.org/wiki/Travelling_salesman_problem}{Travelling Salesman Problem (wikipedia)}
      \end{itemize}
  \end{itemize}
\end{frame}


\begin{frame}{Preguntas}
  \centering
  \Huge
  \textbf{Q\&A}
\end{frame}

\end{document}